\section{Introduction}

The rapid proliferation of digital marketplaces and the growing adoption of the sharing economy have fundamentally reshaped the landscape of personal mobility services worldwide. Traditional vehicle rental models, which rely on centralized intermediaries, fixed-location depots, and manual verification workflows, increasingly fail to meet the expectations of modern consumers who demand convenience, transparency, and trust in peer-to-peer transactions \cite{kauffman2020research}.

This paper presents the design, architectural rationale, and implementation approach of \textbf{LuxeWay} --- a full-stack peer-to-peer (P2P) e-commerce platform for vehicle rental and sharing targeting the Vietnamese market. LuxeWay addresses critical shortcomings identified in existing platforms by combining a hybrid Car Rental and Car Sharing business model with a secure, scalable web architecture, localized payment infrastructure, and a multi-dimensional trust mechanism.

While the theoretical foundations of this work draw upon blockchain, smart contract, and telematics research from the academic literature, the implemented system deliberately applies these concepts through practical, production-grade technologies --- specifically a \textbf{React 18 + TypeScript} frontend and a \textbf{Spring Boot 3 / Java 21} backend backed by \textbf{Microsoft SQL Server} --- in order to deliver a deployable, maintainable platform within the project constraints. This alignment between theoretical motivation and pragmatic engineering decisions is a central contribution of the LuxeWay project.

The remainder of this document is organized as follows. Section~2 defines the system scope. Section~3 analyzes existing platforms and their limitations. Section~4 reviews the relevant academic literature. Section~5 provides a comparative analysis of related research. Section~6 identifies remaining research gaps. Section~7 presents the proposed system architecture and implementation. Section~8 concludes the paper.
